\begin{table}[htbp]
\centering
\footnotesize
\caption{\textbf{6.3} Fracci\'on de clientes de cada puntaje que compr\'o entre el 1 de octubre y el 8 de diciembre. La tasa base es 50.44\,\%.}
\label{tab:p6_fraccion}
\begin{tabular}{ccccc}
\toprule
\textbf{puntaje} & \textbf{R} & \textbf{F} & \textbf{M} & \textbf{$n$ por quintil} \\
\midrule
1 & 27.3\,\% & 31.7\,\% & 28.6\,\% & 732 \\
2 & 37.9\,\% & 31.9\,\% & 38.3\,\% & 731 \\
3 & 54.4\,\% & 44.3\,\% & 44.3\,\% & 732 \\
4 & 65.7\,\% & 58.7\,\% & 60.7\,\% & 731 \\
5 & 66.9\,\% & 85.7\,\% & 80.3\,\% & 732 \\
\bottomrule
\end{tabular}
\\[0.3em]{\scriptsize \textbf{Qu\'e leer:} F es la que m\'as separa: va de 31.7\,\% a 85.7\,\%. R satura arriba (65.7\,\% y 66.9\,\% en los puntajes 4 y 5) pero funciona muy bien abajo: su quintil menos reciente compra con 27.3\,\%, el valor m\'as bajo de la tabla. R sirve para descartar, F para seleccionar.}
\end{table}
